\documentclass[12pt]{article}

\usepackage{amsmath, amssymb}
\usepackage{graphicx}
\usepackage{booktabs}
\usepackage{geometry}

\geometry{
    a4paper,
    margin=1in
}

\title{Proposed Effectiveness Evaluation Framework}
\author{}
\date{}

\begin{document}

\maketitle

\section{Proposed Effectiveness Evaluation Framework}

This work introduces a general framework for evaluating the \textbf{effectiveness}
of population-based metaheuristic algorithms.  
Unlike classical benchmarking, which focuses solely on final objective values,
the proposed framework analyses the \textbf{optimization dynamics}
through normalized diversity, exploration--exploitation behaviour,
AUC indicators, and convergence iteration.

\section{Normalized Positional Diversity}

Let $X=\{x_1,\dots,x_N\}\subset \mathbb{R}^d$ denote a population
and let $\ell, u \in \mathbb{R}^d$ be the lower and upper bounds.  
Each individual is normalized as:
\[
y_i = \frac{x_i - \ell}{u - \ell}.
\]

Pairwise distances are computed as:
\[
d_{ij} = \frac{\|y_i - y_j\|_2}{\sqrt{d}}.
\]

The proposed normalized diversity index is:
\[
D^{*} = 
\frac{2}{N(N-1)} 
\sum_{1 \le i < j \le N} d_{ij}.
\]

Relative diversity progression is defined as:
\[
D_{\text{rel}}(t)
= 
\frac{D^{*}(t)}{D^{*}(0) + \varepsilon},
\]
where $\varepsilon$ prevents division by zero.
A rapid decay of $D_{\text{rel}}(t)$ indicates premature convergence.

\section{Exploration--Exploitation Quantification}

Exploration and exploitation are derived directly from relative diversity:
\[
E(t) = 100 \cdot D_{\text{rel}}(t), 
\qquad 
X(t) = 100 \cdot (1 - D_{\text{rel}}(t)).
\]

High $E(t)$ reflects wide, global search behaviour,
while high $X(t)$ corresponds to intensified local search.

\section{AUC Indicators}

To summarise exploration and exploitation over the entire optimization horizon,
we compute the Area Under Curve (AUC):
\[
\mathrm{AUC}_E 
= \frac{1}{T} \int_0^T E(t)\, dt,
\qquad
\mathrm{AUC}_X 
= \frac{1}{T} \int_0^T X(t)\, dt.
\]

These provide single scalar descriptors of how much exploration
or exploitation occurred throughout the optimization process.

\section{Convergence Iteration}

Convergence is detected when relative diversity falls below a threshold~$\tau$:
\[
D_{\text{rel}}(t) < \tau.
\]
The first iteration that satisfies this condition is recorded as
the \emph{convergence iteration}.

\section{Evaluation Procedure}

Each benchmark function is executed over multiple independent runs.
For every run we record:

\begin{itemize}
    \item best objective value,
    \item final and average diversity,
    \item mean exploration and exploitation,
    \item AUC of exploration and exploitation,
    \item number of function evaluations (FE),
    \item convergence iteration.
\end{itemize}

These aggregated metrics (\textit{mean}, \textit{std}, \textit{median}, \textit{best}, \textit{worst})
yield an overview of algorithmic behaviour
independent of population size, dimensionality, or search-space scale.

\section{Application to the Grey Wolf Optimizer}

When applied to the Grey Wolf Optimizer (GWO) on benchmark function F1,
the proposed framework revealed a very rapid collapse in positional diversity,
followed by an early shift into full exploitation.
This behaviour aligns with the well-known tendency of GWO to experience premature convergence
and highlights the sensitivity and diagnostic value of the proposed evaluation framework.

\end{document}
